\documentclass[a4paper,11pt]{article}
\usepackage[margin=27mm]{geometry}
\usepackage{fontspec,amsmath,amssymb,amsthm,unicode-math,booktabs}
\setmainfont{STIXGeneral}\setmathfont{STIX Two Math}
\newcommand{\gen}[1]{\langle#1\rangle}
\newcommand{\Stab}{\operatorname{Stab}}
\newcommand{\Norm}{\operatorname{N}}
\newcommand{\Aut}{\operatorname{Aut}}
\newtheorem{theorem}{Theorem}\newtheorem{lemma}[theorem]{Lemma}\newtheorem{corollary}[theorem]{Corollary}\newtheorem{proposition}[theorem]{Proposition}
\title{Intersections in simultaneous group actions}
\author{Human Compression project}\date{Author edition; pending independent review}
\begin{document}\maketitle
\begin{abstract}
When a connected component of a product of two permutation actions is
projected onto its first coordinate, an intersection can be repaired
there before its remaining defect is tested in one fiber. This observation
gives the path criterion for mixes of polytopes without a symmetry
assumption. A normalizer hypothesis then identifies the fibers with
actual groups of automorphisms. We separate this algebraic identification
from the further transitivity needed to reduce the number of flags tested.
The resulting applications are the mix criteria of Cunningham and Hubard;
the elementary image-and-kernel argument, rather than a claim of
historical novelty, is the common method.
\end{abstract}

\section{Repairing the first coordinate}
Let a group $W$ act on the left on sets $X$ and $Y$. It acts diagonally on
$X\times Y$. For $S\leq W$ and $q=(x,y)$ write
\[
 S q=\{(s x,s y):s\in S\},\qquad
 D_S(q)=\{y':(x,y')\in S q\}=(S\cap\Stab_W(x))y.
\]
The last equality is literal: a word reaches a point with first coordinate
$x$ precisely when that word fixes $x$. No regularity, finiteness or
normality is assumed.

\begin{theorem}\label{fiber}
Let $A,B,C\leq W$, with $C\leq A\cap B$, and fix $q=(x,y)$.
Suppose that $Ax\cap Bx=Cx$. Then
\[
 Aq\cap Bq=Cq\quad\Longleftrightarrow\quad
 D_A(q)\cap D_B(q)=D_C(q).                         \tag{1}
\]
\end{theorem}
\begin{proof}
The forward implication follows by taking the points with first coordinate
$x$ in the asserted equality. Conversely let $q'=(x',y')\in Aq\cap Bq$.
Its first coordinate belongs to $Ax\cap Bx$, so choose $c\in C$ with
$x'=cx$. Then $c^{-1}q'$ belongs to both $Aq$ and $Bq$, since $c$ lies
in both subgroups. Its first coordinate is $x$, so the equality of the
three fibers gives $c^{-1}q'\in Cq$. Multiplying by $c$ gives $q'\in Cq$.
The opposite inclusion follows from $C\leq A\cap B$.
\end{proof}

This proof is the action version of a familiar subgroup argument. If two
subgroups have the expected common image under a homomorphism, subtract
a lift of that image and test only the remaining kernel. In (1), the
first action replaces the quotient homomorphism. Stabilizers need not be
normal, so the first formulation is in terms of sets rather than quotient
groups.

\section{Recursive path intersections}
Put $I=\{0,\ldots,n-1\}$ and let
\[
 W=W_n=\gen{r_i\ (i\in I)\mid r_i^2=1,\ [r_i,r_j]=1\ (|i-j|>1)}.
\]
For $J\subseteq I$, put $W_J=\gen{r_j:j\in J}$.
A connected premaniplex is a transitive $W$-set. Its colored adjacencies
are the actions of the $r_i$; loops and coincident differently colored
edges are allowed. A maniplex additionally has no fixed point of any
$r_i$ or $r_ir_j$ with $i\ne j$. We retain the standard theorem of
Garza-Vargas and Hubard: a maniplex is the flag graph of a polytope
exactly when its color paths have the intersection property. The interval
form used below is
\[
 W_{[0,j]}x\cap W_{[i,n-1]}x=W_{[i,j]}x.
                                                               \tag{2}
\]
The corresponding cases with a single common color are included. In
particular rank two maniplexes are polygons, finite or infinite, and
satisfy the condition. This geometric characterization is an attributed
input; it is not part of the algebra in Theorem~\ref{fiber}.

\begin{lemma}\label{recursive}
For an $n$-maniplex $Q$, $n\geq3$, the following are equivalent:
\begin{enumerate}
\item $Q$ is polytopal;
\item its facets are polytopal and, at every flag $q$ and every
$0\leq i\leq n-2$,
\[
 W_{[0,n-2]}q\cap W_{[i,n-1]}q=W_{[i,n-2]}q;
\]
\item its facets and vertex-figures are polytopal and, at every flag $q$,
\[
 W_{[0,n-2]}q\cap W_{[1,n-1]}q=W_{[1,n-2]}q.          \tag{3}
\]
\end{enumerate}
\end{lemma}
\begin{proof}
Polytopality implies the other conditions by (2), including its application
inside faces. To prove (2) from the second condition, take a flag $q'$ in
its left-hand side for arbitrary $i,j$. If $j=n-1$ there is nothing to
prove. Otherwise the second condition puts $q'$ in $W_{[i,n-2]}q$.
Inside the common facet, its intersection with $W_{[0,j]}q$ is exactly
$W_{[i,j]}q$, because that facet is polytopal.

The third condition implies the second. The case $i=0$ is immediate.
For $i\geq1$, a flag in
$W_{[0,n-2]}q\cap W_{[i,n-1]}q$ lies, by (3), in
$W_{[1,n-2]}q$. Inside the vertex-figure, intersecting that orbit with
$W_{[i,n-1]}q$ gives $W_{[i,n-2]}q$. This proves all assertions, using
the retained path characterization only at the beginning and end.
\end{proof}

Let $X$ be a polytope, $Y$ a connected premaniplex, and let $Q$ be the
orbit of a chosen pair in $X\times Y$. This is their rooted mix.
It is a maniplex: a fixed point of a generator or a product of two
distinct generators would project to such a fixed point in $X$.
Every restricted-color component of $Q$ is the mix of the corresponding
rooted components of $X$ and $Y$.

\begin{corollary}\label{mix}
Write $A=W_{[0,n-2]}$, $B=W_{[1,n-1]}$, $C=W_{[1,n-2]}$.
Then $Q$ is polytopal if and only if its facets and vertex-figures are
polytopal and
\[
 D_A(q)\cap D_B(q)=D_C(q)\quad\hbox{for every flag }q\in Q. \tag{4}
\]
\end{corollary}
\begin{proof}
The first factor is polytopal, so $Ax\cap Bx=Cx$ at every flag $x$.
Apply Theorem~\ref{fiber} to (3) at every flag of the mix, then apply
Lemma~\ref{recursive}.
\end{proof}
There is no admissibility assumption in this criterion. It keeps the
fiber as a set when there is no appropriate group of automorphisms.

\section{When the fibers are groups}
Fix base flags $x,y$ and put $N_1=\Stab_W(x)$, $N_2=\Stab_W(y)$.
The action on $Y$ is the action on the left cosets $W/N_2$.
If $u\in\Norm_W(N_2)$, the rule
\[
 (gN_2)[u]=guN_2
\]
defines an automorphism commuting with the left $W$-action. Conversely,
an automorphism commuting with $W$ is determined by the image of $N_2$;
that image is $uN_2$ for some $u$ normalizing $N_2$. Thus the group of
these right automorphisms is $\Norm_W(N_2)/N_2$. Its action is free on
flags: one fixed flag determines the identity by transitivity.

Suppose only that $N_1\leq\Norm_W(N_2)$. For $S\leq W$ define the
actual subgroup
\[
 V_S(q)=\{[u]\in\Norm_W(N_2)/N_2:u\in N_1\cap S\}.       \tag{5}
\]
The evaluation $[u]\mapsto y[u]$ identifies $V_S(q)$ bijectively with
$D_S(q)$. This follows directly from the definitions; its injectivity
uses the free action just proved. Consequently
\[
 D_A(q)\cap D_B(q)=D_C(q)
 \quad\Longleftrightarrow\quad V_A(q)\cap V_B(q)=V_C(q).    \tag{6}
\]
Here $V_C\leq V_A\cap V_B$, so inclusion in $V_C$ is an equivalent test.
This also repairs the proof of [CH, Proposition~5.2]. Normality of the
denominator in the product is sufficient: if $u\in N_2$ and $v\in N_1$,
then $uvN_2=vN_2$ since $v$ normalizes $N_2$. No reverse
normalization is inferred from this quotient.
At another flag $gq$, both stabilizers are conjugated by $g$, and the
normalizing hypothesis persists. Formulae (5)--(6) therefore apply
at every flag; they must be formed at that flag, not just at the original
base.

\begin{lemma}\label{sections}
Under the normalizing hypothesis, the lower variance group of the
$S$-components in $Y$ relative to those in $X$ is well-defined and
embeds in $\Aut(Y)$ with image $V_S(q)$.
\end{lemma}
\begin{proof}
The two component stabilizers in $S$ are $N_i\cap S$. Since
$N_1$ normalizes $N_2$, $N_1\cap S$ normalizes $N_2\cap S$.
The lower variance group is
\[
 (N_2\cap S)(N_1\cap S)/(N_2\cap S).
\]
Sending the class of $u\in N_1\cap S$ to its class modulo $N_2$ is
well-defined and injective: its kernel is exactly $N_1\cap N_2\cap S$.
Its image is (5). This proves the actual embedding rather than merely
an abstract isomorphism of the variance groups. In particular the
intersection in (6) takes place inside one specified automorphism group.
\end{proof}

The same quotient calculation supplies the upper variance group. Under
$N_1\leq\Norm_W(N_2)$, the map
\[
 N_1/(N_1\cap N_2)\longrightarrow N_2N_1/N_2,
 \qquad u(N_1\cap N_2)\longmapsto uN_2
\]
is a well-defined surjective homomorphism with trivial kernel. It is
therefore an isomorphism, with no cardinality argument. The group is
trivial exactly when $N_1\leq N_2$, which is exactly a rooted covering
$X\to Y$. If instead $N_2\leq N_1$ and the quotient is defined, both
variance groups are $N_1/N_2$. These statements include the source's
covering examples and Proposition~5.4. For regular inputs the upper
group is the kernel of the projection of their mixed marked groups,
and the lower group is the kernel of the projection to their largest
common quotient $W/(N_1N_2)$.

For finite polygons with $p$ and $q$ sides, the two stabilizers in the
infinite dihedral group are the cyclic groups generated by the $p$th
and $q$th powers of its rotation. Their intersection is generated by the
least-common-multiple power, and their product by the greatest-common-
divisor power. This proves the source's Example~5.6, including lower
variance order $q/\gcd(p,q)$ for $X(M_2\mid M_1)$ when $M_1$ is the
$p$-gon and $M_2$ the $q$-gon. Reversing the arguments gives
$p/\gcd(p,q)$, exactly as printed there. For an infinite polygon its stabilizer is trivial,
so the same assertions are read directly as subgroup statements rather
than as arithmetic with an undefined infinite greatest common divisor.

If $N_1,N_2$ are both normal in a common subgroup $V\leq W$, then
each normalizes the other. The premaniplex $W/V$ is a common symmetry
quotient, and the terminology of Cunningham--Hubard is that both inputs
are $T$-admissible for $T=W/V$. The stabilizer of the mix is
$N=N_1\cap N_2$, and $V/N$ acts by right automorphisms on $W/N$.
The diagonal map $V/N\to V/N_1\times V/N_2$ is injective, since its
kernel is precisely $N$. Each restricted component has the same property
with $V\cap S$ in place of $V$. These facts recover the admissibility
and subgroup identifications used in the source's main criterion.


The restriction statements of [CH, Proposition~2.6] follow as well.
For orientability, a coloring is a character $\epsilon_I$ vanishing on
the flag stabilizer. Restrict that character and stabilizer to $W_J$;
if the character becomes zero, use the one-color coloring. This is an
explicit boundary convention for Proposition~2.6: the restricted quotient
is the connected one-flag component. We do not assert a surjective covering
of two disconnected flags with only semi-edges when the restricted index
set contains every color. For the rotary assertion, a regular component remains regular. Otherwise the
input has two flag orbits of class $2_I$. Its flag stabilizer is normal
in $V=\ker\epsilon_I$, so $N\cap W_J$ is normal in $V\cap W_J$.
The latter subgroup has index at most two in $W_J$. Consequently the
component is regular or has two flag orbits with precisely the restricted
color character. If a larger normalizer merges those two orbits, it is
regular, which is still rotary. Combining these two arguments gives the
admissibility assertion, for facets and vertex-figures with the usual
renumbering of their colors.

\section{How many flags must be tested?}
At this point there are two different reductions: (6) replaces fibers
by groups; a further orbit argument may reduce the number of flags at
which (6) is tested. They have different hypotheses.

\begin{proposition}\label{orbit}
Let $V$ contain and normalize $N_1,N_2$. One representative of each
double coset in $C\backslash W/V$ suffices for the test (4).
In particular one base flag suffices if $W=CV$.
\end{proposition}
\begin{proof}
For $c\in C$, multiplication by $c$ fixes each of the three orbit sets
$Aq,Bq,Cq$ as sets, because $C\leq A\cap B$. Thus equality (3) at
$q$ is equivalent to equality at $cq$. Right automorphisms induced by
$V/N$ preserve all three colored path relations. If $g'=cgv$, then
$g'N=cgvN$ is obtained from $gN$ by such a right automorphism and a
left multiplication by $c$. The condition is therefore constant on each
double coset. Corollary~\ref{mix} proves the assertion.
\end{proof}
The underlying component argument also holds for any set $J$ of colors:
if $W=W_JV$, the right action of $V/N$ is transitive on the
$W_J$-components of the mix. Indeed $g=jv$ places $gN$ in the
component obtained from the base component by $v$. With $J$ the
complement of a chain type, these are the chain components used in
[CH, Proposition~4.6]. This proves its arbitrary-color assertion,
not only the middle-color case needed below.

The condition $W=CV$ says exactly that the common quotient $T=W/V$
is connected on the middle colors. More generally, its $C$-components
index the displayed sufficient list of tests. This assertion does not
infer transitivity of a mix from transitivity of its two factors.

There is a further useful reduction when the mix has automorphisms
which do not project to the two factors. Condition (3) itself is
invariant under every automorphism of the mix and under middle-color
paths. Thus one base flag suffices whenever the mix is medial-section-
transitive. In particular it suffices when the mix is regular, regardless
of whether its full automorphism group acts coordinatewise. This uses
only the path equality, followed by its already proved equivalence
with the fiber or variance test at that base.

For regular inputs take $V=W$, so there is one test. For two-orbit
inputs in a fixed class, take $V$ to be the kernel of the character
$\epsilon_I:W\to C_2$ with $\epsilon_I(r_i)=0$ for $i\in I$ and
$1$ otherwise. If some middle color is outside $I$, then $CV=W$.
If all middle colors lie in $I$, there are two tests, represented by
$1$ and any $r_j$ with $j\notin I$. The possible complements of $I$
are then $\{0\},\{n-1\},\{0,n-1\}$. Applying the same argument
with $A$ and $B$ describes whether there is one orbit or two of facets
and vertex-figures. This gives the three exceptional tables in the
application without silently identifying different rooted sections.

\section{A useful quotient consequence}
The variance group in (5) is trivial exactly when every word of $S$
fixing the first base flag also fixes the second. Equivalently the
first rooted $S$-component covers the second. Thus a covering relation
between facets makes the facet variance trivial.

\begin{proposition}\label{facetcover}
Suppose the two inputs are admissible over a common $T$, the first is
polytopal, and at every paired flag its first facet covers its second
facet. Then their mix is polytopal.
\end{proposition}
\begin{proof}
Proceed by induction on the rank, starting with polygons. The mix of
a polygon and a rank-two premaniplex is a rank-two maniplex, hence a
polygon. In higher rank the mixed facets are the first facets and are
polytopal. The facets of each mixed vertex-figure are mixes of paired
medial sections. The assumed facet coverings restrict to coverings
between these sections. Admissibility restricts to the vertex-figures,
so induction shows that all mixed vertex-figures are polytopal.
Finally $V_A(q)=1$ at every flag, so (6) and Corollary~\ref{mix} give
polytopality of the whole mix. Interchanging the two ends proves the
dual result.
\end{proof}

\section{Attribution and remaining application checks}
The path characterization is retained from Garza-Vargas--Hubard.
In the source's symmetry-type discussion on physical page~4, adjacency
between two flag orbits means that each flag has its unique neighbor
in the other orbit, not that it is adjacent to every flag there. This
is immediate because automorphisms commute with the color involutions.
The mix, admissibility and variance terminology, the recursive criteria
and the principal polytopality theorem are due to Cunningham--Hubard.
The first-coordinate proof above is their correction-of-a-path argument
expressed for arbitrary group actions; it is also the same elementary
image-and-kernel mechanism used in our supported permutation proofs.
The one-sided normalizer formulation states exactly the algebra needed.
We make no historical novelty claim for it.

The companion application binds the printed single-base corollaries to
a justified orbit hypothesis and exhibits their failing literal special
cases. Its source table separately identifies the criteria in which
actual roots replace suppressed identifications. Nothing here proves any of the source's three
concluding open problems. This working edition carries no source count,
conjecture count or numerical comparison.

\begin{thebibliography}{9}
\bibitem{CH} G. Cunningham and I. Hubard, \emph{Polytopality criteria for the
mix of polytopes and maniplexes}, arXiv:2506.19334v1, 24 June 2025.
\bibitem{GH} J. Garza-Vargas and I. Hubard, \emph{Polytopality of maniplexes},
Discrete Math. 341 (2018), 2068--2079, Theorem 5.3; retained through
[CH, Theorem 2.1].
\end{thebibliography}
\end{document}
